\section{A small polynomial-weight parity network}\label{sec:paritycircuit}
For completeness, full parity on $k$ bits has a one-hidden-layer
$\tanh$ realization of width $k+1$, with weights $O(k\log(k+1))$
and output margin greater than $1/2$. Let
\[
 t(x)=\frac{k+\sum_i x_i}{2}\in\{0,\ldots,k\},\quad
 y_j=(-1)^{k-j},\quad A=\log(16(k+1)).
\]
The step functions $s_j(t)=\sign(t-j+1/2)$ satisfy
\[
 y_t=c+\sum_{j=1}^k y_js_j(t),\qquad c=(y_0+y_k)/2\in\{0,1\}.
\]
Replace each step by $h_j=\tanh(A(t-j+1/2))$. Since its argument
has absolute value at least $A/2$, its error is at most $2e^{-A}$.
A constant neuron $h_0=\tanh1$ supplies the term $c$ using output
weight $c/\tanh1$. The resulting output differs from parity by
less than $1/8$. First-layer input weights are $A/2$ and bias
weights have magnitude at most $Ak/2$; readout weights have
magnitude at most $2$.

At depth $L\ge3$, the next hidden layer sums these features and
applies $\tanh$, producing the correct sign with magnitude at least
$\tanh(7/8)>1/2$. Each additional hidden layer can preserve this
margin by a single weight $2$, since $\tanh(2u)>1/2$ whenever
$u>1/2$. A final readout weight $1$ completes the construction.
Unused neurons are set to zero. Thus the same polynomial bounds
hold at every fixed depth.
For any other fixed margin $a>0$, multiplying the final readout by $2a$
gives margin greater than $a$ with the same polynomial complexity bounds.

